\section{Meyer's quotient, its jets and its receivers}\label{sec:quotient}

Throughout this section $\B$, $I_D$, $I_\zeta$ and $Q=\B/I_\zeta$ are as in Section~\ref{sec:setting}, and $L$ denotes multiplication by $s$ on a quotient $\B/I_D$.
Since the seminorms $b_{A,M}$ increase in $A$ and $M$, they form a directed family, and every quotient of $\B$ by a closed subspace is a Fréchet space with the quotient seminorms $\bar b_{A,M}$ (Rudin \cite{Rudin}, Theorem~1.41).

\subsection{The strip space}

\begin{lemma}[$\B$ is a Fréchet--Schwartz space]\label{lem:schwartz}
For all $A,M\ge0$ the unit ball of $b_{A+1,M+1}$ is precompact for $b_{A,M}$.
Consequently $\B$ and every quotient $\B/I$ by a closed subspace are Fréchet--Montel spaces, hence reflexive, and $(\B/I)'_\beta\cong I^\perp\subset\B'_\beta$ topologically.
\end{lemma}
\status{Proved here (claude-ab, \note{27\_} Lemma 27.2; refereed in the ninth pass). The programme proves the same properties for $\mathcal{A}$ and transfers them to $\B/I$ through the summation image (\lbl{SDT}); the lemma removes that dependency. The consequences are standard permanence properties of Schwartz spaces.}

\begin{proof}
Let $b_{A+1,M+1}(F)\le1$, so $|F(s)|\le(1+|t|)^{-M-1}$ on $|\re s|\le A+1$.
Cauchy's estimate on discs of radius $\tfrac12$ gives $|F'(s)|\le2(\tfrac12+|t|)^{-M-1}$ on $|\re s|\le A+\tfrac12$, so the ball is bounded and equicontinuous there.
Given $\varepsilon>0$ choose $T$ with $(1+T)^{-1}<\varepsilon$; then $(1+|t|)^M|F|<\varepsilon$ on $|t|\ge T$ for every $F$ in the ball.
On the rectangle $|\re s|\le A$, $|t|\le T$, Arzelà--Ascoli gives a finite $\varepsilon(1+T)^{-M}$-net for the supremum norm; with the tail estimate it is a $2\varepsilon$-net for $b_{A,M}$.
A Fréchet space in which every seminorm ball is precompact for the preceding seminorm is a Schwartz space; Schwartz spaces are Montel, and quotients of Fréchet--Schwartz spaces by closed subspaces are Fréchet--Schwartz (Jarchow \cite{Jarchow}); Montel spaces are reflexive (Schaefer \cite{Schaefer}, Ch.~IV).
Bounded sets of $\B/I$ are relatively compact, and compact sets lift through the quotient map of a Fréchet space, so the strong topology of $(\B/I)'$ is the one induced from $\B'_\beta$ on $I^\perp$.
\end{proof}

\begin{proposition}[multiplication by $s$ on $\B/I_D$]\label{prop:spectral-L}
Let $\Phi\in\B$ have zero divisor exactly $D$, and let $L$ be multiplication by $s$ on $Q_D=\B/I_D$.
\begin{enumerate}[label=(\alph*)]
\item For $\lambda\notin\operatorname{supp}D$, $\lambda-L$ is invertible on $Q_D$, with inverse induced by $R_\lambda F=(F-\Phi\cdot F(\lambda)/\Phi(\lambda))/(\lambda-s)$.
\item Every $\rho\in\operatorname{supp}D$ is an eigenvalue; hence the spectrum of $L$ is $\operatorname{supp}D$.
\item If $\rho$ has multiplicity $m$ in $D$, then $\ker(L-\rho)^m=\ker(L-\rho)^{m+1}$ has dimension $m$, the jet map $j_\rho$ identifies it with $\C[T]/(T^m)$, and $L-\rho$ acts on it as multiplication by $T$, a single nilpotent Jordan block of length $m$.
\end{enumerate}
\end{proposition}
\status{Proved in the programme for the zeta divisor (\lbl{RZ3--RZ7}, including the Riesz projector at $\rho$ as multiplication by an entire function $e_\rho$; audited in \note{27\_}); the general-divisor form and the proof given here are claude-ab's (\note{24\_} Lemma 24.2 step 1, \note{27\_} Lemma 27.1).}

\begin{proof}
(a) The numerator of $R_\lambda F$ vanishes at $\lambda$, so $R_\lambda F$ is entire; it lies in $\B$ (divide by $|\lambda-s|\ge1$ away from $\lambda$, and use the maximum principle on $|s-\lambda|\le1$), and it preserves $I_D$ because $\Phi\in I_D$ and $\lambda\notin\operatorname{supp}D$.
Since $(\lambda-s)R_\lambda F=F-\Phi F(\lambda)/\Phi(\lambda)\equiv F$ and $R_\lambda((\lambda-s)F)=F$, $R_\lambda$ induces the inverse of $\lambda-L$.
(b) $U=\Phi/(s-\rho)\in\B$ vanishes to order $m-1$ at $\rho$, so $U\notin I_D$, while $(s-\rho)U=\Phi\in I_D$.
(c) A class $[F]$ lies in $\ker(L-\rho)^k$ exactly when $(s-\rho)^kF$ vanishes to full order on $D$, that is, when $F$ vanishes to full order at every other point of $D$ and to order $m-k$ at $\rho$. So $\ker(L-\rho)^m=\ker(L-\rho)^{m+1}$ consists of the classes of functions vanishing to full order off $\rho$, and $j_\rho$ maps it injectively into $\C[T]/(T^m)$.
It is onto: $U_\rho=\Phi/(s-\rho)^m\in\B$ has $U_\rho(\rho)\ne0$ and vanishes to full order at the other points, so $U_\rho\cdot P(s-\rho)$ with $P=[P_0\,u(T)^{-1}]_{<m}$, $u$ the jet of $U_\rho$, has jet $P_0$ at $\rho$.
Finally $j_\rho((s-\rho)F)=T\,j_\rho(F)$.
\end{proof}

\begin{lemma}[no equivariant maps across disjoint divisors]\label{lem:disjoint}
Let $D$ be the divisor of some $F_D\in\B$ and let $D'$ be a divisor with $D\cap D'=\emptyset$.
Every linear map $H:\B/I_D\to\B/I_{D'}$ with $HL=LH$ is zero. Continuity is not needed.
\end{lemma}
\status{Proved here (claude-ab, \note{24\_} Lemma 24.2; this is the uniqueness argument of the programme's \lbl{GSL6}, which it strengthens: only commutation with $L$ is used). Refereed (eighth pass).}

\begin{proof}
Let $\lambda\in D'$ have multiplicity $m$. By Proposition~\ref{prop:spectral-L}(a), $(L-\lambda)^m$ is onto $\B/I_D$. For every $x$, $(L-\lambda)^mH(x)=H((L-\lambda)^mx)$ has zero $m$-jet at $\lambda$; since $x$ ranges over the image of $(L-\lambda)^m$, every $H(y)$ has zero $m$-jet at every $\lambda\in D'$, so $H(y)\in I_{D'}$.
\end{proof}

\begin{lemma}[uniqueness of exponential sums on a divisor]\label{lem:expsum}
Let $0\not\equiv G\in\B$, let $Z$ be a set of zeros of $G$ in a strip $|\re s|\le A$, and let $(c_\rho)\in\ell^1(Z)$.
If $\sum_\rho c_\rho e^{\rho x}=0$ for every real $x$, then $c=0$.
\end{lemma}
\status{Proved in the programme for the Gaussian-weighted zeta divisor (\lbl{HBW8.9--8.12}; \lbl{HSR4.1} reduces to it); stated for general $G$ and proved here by claude-ab (\note{32\_} Lemma 32.2). Refereed (eleventh pass).}

\begin{proof}
Every $F\in\B$ has a half-Mellin representation $F(s)=\tfrac12\int_{\R} a(e^v)e^{sv}\,dv$ with $a(e^v)$ decaying faster than every exponential (\lbl{NJS4.1}). By the $\ell^1$ condition and $|e^{\rho v}|\le e^{A|v|}$ the interchange is justified and $\sum_\rho c_\rho F(\rho)=\tfrac12\int a(e^v)\sum_\rho c_\rho e^{\rho v}\,dv=0$.
For $\rho_0\in Z$ of order $m$ as a zero of $G$, the function $F=G/(s-\rho_0)^m\in\B$ has $F(\rho_0)\ne0$ and vanishes on $Z\setminus\{\rho_0\}$, so $c_{\rho_0}=0$.
\end{proof}

\subsection{The exact summation image}

Let $S$ be the space of even Schwartz functions $f$ on $\R$ with $f(0)=0$ and $\int f=0$, put $\Sigma f(u)=2\sum_{n\ge1}f(nu)$, and let $\mathcal{M}_0b(s)=\int_0^\infty b(u)u^{s-1}du$.

\begin{proposition}[the exact summation image]\label{prop:summation-image}
$f\mapsto\mathcal{M}_0\Sigma f$, which equals $2\zeta(s)\int_0^\infty f(v)v^{s-1}dv$ on $\re s>1$, is a topological isomorphism from $S$ onto $I_\zeta$. Its inverse is $F\mapsto f_F$ with
\[
f_F(x)=\frac1{2\pi}\int_{\R}\frac{F(2+it)}{2\zeta(2+it)}x^{-2-it}\,dt=\frac12\sum_{n\ge1}\mu(n)\,(\mathcal{M}_0^{-1}F)(nx)\quad(x>0),
\]
and $f_F^{(2r)}(0)/(2r)!=F(-2r)/(2\zeta'(-2r))$ for $r\ge1$, $\int_0^\infty f_F\,dx/x=-F(0)$, $\int_0^\infty f_F\log x\,dx=F(1)/2$.
In particular the summation image is closed: there is no gap between it and its closure.
\end{proposition}
\status{Proved in the programme (\lbl{SSI1--SSI7}, \lbl{OMS2--OMS4}); audited (\note{25\_} Lemma 25.1, \note{26\_}); the Taylor coefficients and endpoint identities checked on the programme's source vector. On the programme's own account (\lbl{OMS0}) this is Meyer's range theorem \cite{Meyer} made explicit for $\B$, with the division estimate supplied; the audit did not read Meyer's paper. Refereed (eighth pass).}

\emph{Outline of the proof.}
Poisson summation, $\Sigma\hat f(u)=u^{-1}\Sigma f(u^{-1})+u^{-1}f(0)-\int f$, controls $u\to0$, so $\Sigma$ maps $S$ continuously into the functions whose Mellin transforms lie in $\B$; on $\re s>1$ the transform is Müntz's $2\zeta\mathcal{M}f$ (Titchmarsh \cite{Titchmarsh}, §2.11), with $\mathcal{M}f$ holomorphic on $\re s>-2$ and zero at $1$, so it vanishes to full order on $\Zs$, and the Euler product gives injectivity.
Conversely, for $F\in I_\zeta$ the quotient $H=F/(2\zeta)$ is holomorphic except for simple poles at $-2r$; $1/\zeta$ is polynomially bounded on the two edges of $-2N-1\le\re s\le b$ (by the reflection formula for $|\chi(-2N-1+it)|$ and $|1/\zeta(b+it)|\le\zeta(b)$), and inside the strip the division estimate $|F/F_*|\le C\exp(C(|t|+2)^{3/2}\log(|t|+2))$ (\lbl{S3}; proved by the Hadamard product and a minimum-modulus argument, \lbl{OMS2.8}) is removed by the factor $e^{\varepsilon s^2}$ and the maximum principle, giving a bound linear in a seminorm of $F$.
Shifting the inverse Mellin contour to $\re s=-2N-1$ picks up the residues $F(-2r)/(2\zeta'(-2r))\,x^{2r}$, so $f_F$ extends to an even Schwartz function with $f_F(0)=0$ and $\int f_F=2H(1)=0$. \qed

Two consequences are used below.
The closure of Connes' summation map in the rapid-decay Fréchet space is the ideal of all zero-jets (\lbl{S5} of the second attempt; register S13), and by Proposition~\ref{prop:summation-image} the image itself is already that ideal.
And Meyer's quotient $Q=\mathcal{A}/\mathcal{E}S$ is identified with $\B/I_\zeta$, on which the jet map is injective by the definition of $I_\zeta$.

\subsection{The jet map and the primary decomposition}

For each zero $\rho$ the programme's section $s_\rho:A_\rho\to Q$ (\lbl{OMS6}) is $s_\rho(P)=[U_\rho\cdot(P\,c_\rho)(s-\rho)]$ with $U_\rho=F_*/(s-\rho)^{m_\rho}$ and $c_\rho=u(T)^{-1}\bmod T^{m_\rho}$, $u$ the jet of $U_\rho$ at $\rho$.
It satisfies $j_\eta s_\rho=\delta_{\eta\rho}\operatorname{id}$, so the $\Pi_\rho=s_\rho j_\rho$ are commuting idempotents with $\Pi_\rho\Pi_\eta=0$ for $\rho\ne\eta$, and finite sums of sections realise every finitely supported jet tuple.
On the range of $\Pi_\rho$ the dilations act by $W_a=a^\rho\exp((\log a)N_\rho)$ with $N_\rho=L-\rho$ (Proposition~\ref{prop:spectral-L}(c) and \lbl{OMS6.5}).

\begin{lemma}[jet decay]\label{lem:jetdecay}
For $F\in\B$, a nontrivial zero $\rho=\beta+i\gamma$ and $r,M\ge0$,
\(
(1+|\gamma|)^M\,|F^{(r)}(\rho)/r!|\le2^{r+M}\,b_{2,M}(F).
\)
\end{lemma}
\status{Proved here (claude-ab, \note{40\_} Lemma 40.1; an explicit form of \lbl{OMS5.6}). Refereed (seventeenth pass).}

\begin{proof}
Cauchy's formula on $|z-\rho|=\tfrac12$ gives $|F^{(r)}(\rho)/r!|\le2^r\max_{|z-\rho|=1/2}|F|$. The circle lies in $-\tfrac12<\re z<\tfrac32$ because $0<\beta<1$, and on it $1+|\im z|\ge(1+|\gamma|)/2$, so $|F(z)|\le2^Mb_{2,M}(F)(1+|\gamma|)^{-M}$.
\end{proof}

\begin{corollary}\label{cor:jetimage}
The image $J(Q)$ of the jet map $J=(j_\rho)_\rho:Q\to\prod_\rho A_\rho$ is a proper dense subspace of the product, and the quotient topology of $Q$ is strictly finer than the product topology on $J(Q)$. In particular $Q$ is not the product of its primary blocks.
\end{corollary}
\status{Proved in the programme (\lbl{OMS5}; audited in \note{40\_} §1, no errors); the explicit constant of Lemma~\ref{lem:jetdecay} is claude-ab's.}

\begin{proof}
Density: finite sums of sections realise every finitely supported tuple. Properness: if $F(\rho)=1$ at every zero, Lemma~\ref{lem:jetdecay} with $r=0$, $M=1$ gives $1+|\gamma|\le2b_{2,1}(F)$ for all ordinates, which are unbounded.
Each $j_\rho$ is continuous and vanishes on $I_\zeta$, so the quotient topology is at least as fine as the product topology; if they were equal, $J(Q)$ would be complete, hence closed in the Hausdorff product, hence (being dense) the whole product.
\end{proof}

Two classical facts are used (Titchmarsh \cite{Titchmarsh}, Theorems 9.2 and 9.6(A)):
(D1) the number of zeros, with multiplicity, with $|\gamma-T|<1$ is $O(\log T)$;
(D2) for $-1\le\sigma\le2$ and $t\ge2$, $\zeta'/\zeta(s)=\sum_{|t-\gamma|<1}(s-\rho)^{-1}+O(\log t)$.

\begin{lemma}[good heights; Titchmarsh, Theorem 9.7]\label{lem:goodheights}
There are $A,T_0>0$ such that every interval $[T,T+1]$ with $T\ge T_0$ contains some $t$ with $|\zeta(\sigma\pm it)|\ge t^{-A}$ for all $-1\le\sigma\le2$.
\end{lemma}
\status{Titchmarsh's Theorem 9.7 \cite{Titchmarsh}, with his proof (\note{40\_} Lemma 40.2; the statements of Theorems 9.2, 9.6(A) and 9.7 were checked by a referee on the printed pages).}

\begin{proof}
Integrating the real part of (D2) along $[\sigma+it,2+it]$, for $t$ not an ordinate,
$\log|\zeta(\sigma+it)|=\log|\zeta(2+it)|-\sum_{|t-\gamma|<1}(\log|2+it-\rho|-\log|\sigma+it-\rho|)+O(\log t)$.
Here $|\zeta(2+it)|\ge1/\zeta(2)$, $1<|2+it-\rho|<3$, the number of terms is $O(\log t)$ and $|\sigma+it-\rho|\ge|t-\gamma|$, so $\log|\zeta(\sigma+it)|\ge\sum_{|t-\gamma|<1}\log|t-\gamma|-C\log T$ uniformly in $\sigma$.
Since $\int_{-1}^1\log|u|\,du=-2$, the mean of the sum over $t\in[T,T+1]$ is at least $-2\#\{T-1<\gamma<T+2\}\ge-C'\log T$ by (D1); some $t$, not an ordinate, is at least the mean. For $-t$ use $|\zeta(\bar s)|=|\zeta(s)|$.
\end{proof}

For each zero $\rho$ of multiplicity $m=m_\rho$ let $c_{\rho,k}$ ($0\le k<m$) be the Taylor coefficients at $\rho$ of $(s-\rho)^m/\zeta(s)$; so $\sum_{k<m}c_{\rho,k}(s-\rho)^{k-m}$ is the principal part of $1/\zeta$ at $\rho$, and $c_{\rho,0}=m!/\zeta^{(m)}(\rho)$.
Let $\Lambda_m$ be the Fréchet space of tuples $(P_\rho)\in\prod_\rho A_\rho$ with $\sup_\rho(1+|\gamma|)^M\|P_\rho\|<\infty$ for all $M$, $\|\cdot\|$ the largest coefficient.

\begin{theorem}[when the primary decomposition converges]\label{thm:primary}
The following are equivalent.
\begin{enumerate}[label=(\alph*)]
\item For some enumeration $(\rho_n)$ of the distinct zeros, the partial sums $\sum_{k\le n}\Pi_{\rho_k}x$ converge for every $x\in Q$ (the limit is then $x$).
\item[\textit{(a$'$)}] For every $x\in Q$ the set $\{\Pi_\rho x\}_\rho$ is bounded in $Q$.
\item $\sum_\rho\Pi_\rho x$ converges absolutely for every $x\in Q$.
\item There are $C,K$ with $|c_{\rho,k}|\le C(1+|\gamma|)^K$ for all $\rho$ and all $k<m_\rho$: the principal parts of $1/\zeta$ at the zeros are polynomially bounded.
\item $J:Q\to\Lambda_m$ is an isomorphism of Fréchet spaces.
\end{enumerate}
If all zeros are simple, (c) reads $|\zeta'(\rho)|\ge c(1+|\gamma|)^{-K}$.
In every case (a$'$) implies $|\zeta^{(m_\rho)}(\rho)/m_\rho!|\ge c(1+|\gamma|)^{-K}$.
\end{theorem}
\status{Proved here (claude-ab, \note{40\_} Propositions 40.4 and 40.5; the form for all multiplicities is the seventeenth referee's sharpening, whose proof was checked). Checks F1--F9, G1--G3. Novelty: not found in the programme or in the literature searched (Berenstein--Taylor \cite{BerensteinTaylor}, Abanin--Le Hai Khoi--Nalbandyan \cite{AbaninKhoiNalbandyan}, Bondarenko--Radchenko--Seip \cite{BondarenkoRadchenkoSeip}); it is an instance for Meyer's quotient of the classical link between lower bounds at the zeros and interpolation or representing systems.}

\begin{proof}
\emph{(a)$\Rightarrow$(a$'$) and (b)$\Rightarrow$(a).} $\Pi_{\rho_n}x$ is the difference of consecutive partial sums, and absolute convergence in a Fréchet space implies convergence. If the partial sums tend to $y$, then $j_\rho y=j_\rho x$ for every $\rho$ by continuity, so $y=x$.

\emph{(a$'$)$\Rightarrow$(c).}
Each $\Pi_\rho$ is continuous (Lemma~\ref{lem:jetdecay}), and the family is pointwise bounded on the Fréchet space $Q$, so it is equicontinuous (Banach--Steinhaus; Rudin \cite{Rudin}, Theorem~2.6): there are $A_1,M_1,C_0$ with $\bar b_{2,0}(\Pi_\rho x)\le C_0\bar b_{A_1,M_1}(x)$ for all $\rho,x$.
Take $E_\rho(s)=e^{(s-\rho)^2}$ (the programme's Gaussian test vector, \lbl{CFP6.1}). Since $|E_\rho(\sigma+it)|=e^{(\sigma-\beta)^2-(t-\gamma)^2}$ and $1+|t|\le(1+|\gamma|)(1+|t-\gamma|)$, $b_{A_1,M_1}(E_\rho)\le C_1(1+|\gamma|)^{M_1}$.
So $\Pi_\rho[E_\rho]=s_\rho(e^{T^2}\bmod T^m)$ has a representative $F_\rho$ with $b_{2,0}(F_\rho)\le2C_0C_1(1+|\gamma|)^{M_1}$; it has jet $e^{T^2}$ at $\rho$ and vanishes to full order at every other zero.

Let $\gamma\ge T_0+2$ (negative $\gamma$ are symmetric). Choose good heights $T_1\in[\gamma-2,\gamma-1]$ and $T_2\in[\gamma+1,\gamma+2]$ (Lemma~\ref{lem:goodheights}) and let $R=[-1,2]\times[T_1,T_2]$.
The function $H=F_\rho(s-\rho)^m/\zeta$ is holomorphic on a neighbourhood of $R$: $F_\rho$ cancels the other zeros in $R$, the singularity at $\rho$ is removable, $\zeta\ne0$ on $\partial R$ (the Euler product on $\sigma=2$, the functional equation on $\sigma=-1$, Lemma~\ref{lem:goodheights} on the horizontal sides), and $R$ contains no trivial zero and not the pole.
On $\partial R$: $|F_\rho|\le b_{2,0}(F_\rho)$ and $|s-\rho|<3$; $|1/\zeta|\le\sum n^{-2}=\pi^2/6$ on $\sigma=2$; $|1/\zeta|\le T_2^A$ on the horizontal sides; and on $\sigma=-1$, with $y=t/2$,
\begin{gather*}
|\chi(-1+it)|^2=\pi^{-3}\,y\coth(\pi y)\,(\tfrac14+y^2)\ge\frac1{4\pi^4},\\
|\zeta(-1+it)|=|\chi(-1+it)|\,|\zeta(2-it)|\ge\frac1{2\pi^2}\cdot\frac6{\pi^2}=\frac3{\pi^4},
\end{gather*}
from $\chi(s)=\pi^{s-1/2}\Gamma(\tfrac{1-s}2)/\Gamma(\tfrac s2)$, $|\Gamma(1-iy)|^2=\pi y/\sinh\pi y$ and $|\Gamma(-\tfrac12+iy)|^2=\pi/((\tfrac14+y^2)\cosh\pi y)$.
So $|H|\le B_\rho:=2C_0C_1(1+|\gamma|)^{M_1}3^m\max(\pi^2/6,\pi^4/3,(\gamma+2)^A)$ on $\partial R$, hence on $R$.
The disc $|s-\rho|\le\tfrac12$ lies in $R$, so $|H^{(k)}(\rho)/k!|\le2^kB_\rho$. The jet of $H$ at $\rho$ is $e^{T^2}\sum_kc_{\rho,k}T^k$, and the coefficients of $e^{-T^2}$ have modulus at most $1$, so $|c_{\rho,k}|\le2^{k+1}B_\rho\le2^mB_\rho$.
By (D1), $m_\rho\le C\log(2+|\gamma|)$, so $2^m3^m$ is polynomial in $|\gamma|$; finitely many low zeros are absorbed into the constant, since $c_{\rho,k}$ is finite. The last assertion of the theorem is the case $k=0$.

\emph{(c)$\Rightarrow$(d) and (b).}
$V_\rho=(s-1)\zeta(s)/(s-\rho)^m$ is entire, and the reciprocal of its jet is $v_\rho^{-1}=(\sum_kc_{\rho,k}T^k)(\rho-1+T)^{-1}$, whose coefficients are at most $mC(1+|\gamma|)^K$ because $|\rho-1|\ge1$.
With $\lambda=m+1$ put $q_{\rho,k}=[T^ke^{-\lambda T^2}v_\rho^{-1}]_{<m}$ and
\[
G_{\rho,k}(s)=e^{\lambda(s-\rho)^2}\,V_\rho(s)\,q_{\rho,k}(s-\rho).
\]
$G_{\rho,k}$ is entire, has jet $T^k\bmod T^m$ at $\rho$ and vanishes to full order at every other zero. The coefficients of $q_{\rho,k}$ are at most $Q_\rho=m^2e^{m+1}C(1+|\gamma|)^K$.
For $|\sigma|\le A$, $|s-\rho|\ge1$ and $u=t-\gamma$: $|(s-1)\zeta(s)|\le C_A(1+|t|)^{c_A}$ (polynomial growth in strips; Titchmarsh \cite{Titchmarsh}, Ch.~V), $|q_{\rho,k}(s-\rho)|\le mQ_\rho(A+1+|u|)^{m-1}$, $|e^{\lambda(s-\rho)^2}|\le e^{(m+1)(A+1)^2}e^{-(m-1)u^2}e^{-2u^2}$ and $e^{-u^2}(A+1+|u|)\le A+2$, so
\[
|G_{\rho,k}(s)|(1+|t|)^M\le mQ_\rho\,e^{(m+1)(A+1)^2}(A+2)^{m-1}C_A(1+|\gamma|)^{M+c_A}\sup_ue^{-2u^2}(1+|u|)^{M+c_A};
\]
on $|s-\rho|<1$ the maximum principle applies. Since $m\le C\log(2+|\gamma|)$, $b_{A,M}(G_{\rho,k})\le C_{A,M}(2+|\gamma|)^{N_{A,M}}$.
As $N(T)=O(T\log T)$, $\sum_\rho m_\rho(2+|\gamma|)^{-2}<\infty$, so for $P\in\Lambda_m$ the series $F_P=\sum_\rho\sum_{k<m_\rho}P_{\rho,k}G_{\rho,k}$ converges absolutely in $\B$ and $J[F_P]=P$.
$J$ is continuous into $\Lambda_m$ by Lemma~\ref{lem:jetdecay} (the factor $2^{m_\rho}$ is absorbed by raising $M$) and injective, so it is an isomorphism with inverse $P\mapsto[F_P]$.
Finally $\sum_kx_{\rho,k}[G_{\rho,k}]=\Pi_\rho x$, since both have the same jets at every zero; so $x=[F_{Jx}]=\sum_\rho\Pi_\rho x$, absolutely.

\emph{(d)$\Rightarrow$(b).} $\sum_\rho\sum_kx_{\rho,k}\delta_{\rho,k}$ converges absolutely in $\Lambda_m$, and $J^{-1}$ is continuous with $J^{-1}(\sum_kx_{\rho,k}\delta_{\rho,k})=\Pi_\rho x$.
\end{proof}

\begin{remark}\label{rem:primary}
Whether (c) holds is not known. For simple zeros it reads $|\zeta'(\rho)|\ge c(1+|\gamma|)^{-K}$, and a bound of this form at every zero implies that every zero is simple, which is open; unconditionally, at least a proportion $0.4075$ of the zeros are known to be simple ($0.4058$ by Bui--Conrey--Young \cite{BCY}, as cited in Bui--Heath-Brown \cite{BuiHeathBrown}, §1; $0.407511$ for simple zeros on the line by Pratt--Robles--Zaharescu--Zeindler \cite{PRZZ}). Under RH it is expected that $|\zeta'(\rho)|^{-1}\gg|\gamma|^{1/3-\varepsilon}$ for infinitely many zeros (Milinovich--Ng \cite{MilinovichNg}, p.~2), so such a bound could only hold with $K\ge1/3$.
Whether the leading coefficients ($k=0$) alone suffice when zeros are multiple is not settled; in the weighted spaces $A_p$ the analogous leading-coefficient condition suffices for interpolation (Berenstein--Taylor \cite{BerensteinTaylor}, Theorem~4).
The expansion $x=\sum_\rho\Pi_\rho x$ is formally like the residue expansion of Mertens' function, whose residue at a simple zero, $x^\rho/(\rho\zeta'(\rho))$, carries the same factor $1/\zeta'(\rho)$.
The programme itself uses only finite sums of primary projectors (\lbl{OMS6}, \lbl{SCL7}) and states that it assumes no convergence of an infinite sum (\lbl{GEX2}).
\end{remark}

\subsection{Dense families of test functions}

For fixed $t>0$ put $g_t(s)=e^{ts^2}$ and, for a set $S\subset\{2,3,\dots\}$, $\mathcal{N}_S=\operatorname{span}\{g_t(s)(n-n^{1-s})/s:\ n\in S\}$.

\begin{proposition}[Nyman--Beurling in the Fréchet topology]\label{prop:NB}
If $\limsup_{R\to\infty}\#\{n\in S:\log n\le R\}/R^2=\infty$ (for example $S=\{2,3,\dots\}$, or the $d$-th powers $k^d$, $k\ge2$), then $\mathcal{N}_S$ is dense in $\B$ and $\xi\mathcal{N}_S$ is dense in $I_\zeta$. Moreover $\overline{\xi\B}=I_\zeta$, although $\xi\B\subsetneq I_\zeta$.
\end{proposition}
\status{Proved in the programme (\lbl{NCI2--NCI8}, \lbl{RSS1--RSS2}); audited (\note{16\_}). The sparse version in $\B$ is claude-ab's transplant of the \lbl{NCI8} Jensen step into \lbl{RSS1}. Refereed (fourth pass). The corresponding statement in Noor's Hardy-space setting is equivalent to RH \cite{Noor}; here it is unconditional, because the topology is different.}

\emph{Outline.} A functional $\Lambda$ annihilating the family gives the entire function $L(w)=\Lambda(g_t(e^w-e^{(1-s)w})/s)$ with $\log|L(w)|\le C(1+|w|^2)$, vanishing at $\log n$, $n\in S$. By Jensen's formula a nonzero such function has $O(R^2)$ zeros in the disc of radius $R$, so the growth condition on $S$ forces $L\equiv0$. The identity $(\partial_w-1)(e^w-e^{(1-s)w})/s=e^{(1-s)w}$ then gives $\Lambda(g_te^{vs})=0$ for all $v$, and the half-Mellin representation extends this to $\Lambda=0$. \qed

\begin{proposition}[an infinite set of degrees need not suffice]\label{prop:NBfail}
If $S\subset\{a^k:k\ge1\}\cup F$ with $a\ge2$ and $F$ finite, then $\mathcal{N}_S$ is not dense in $\B$ and $\xi\mathcal{N}_S$ is not dense in $I_\zeta$.
\end{proposition}
\status{Proved here (claude-ab, \note{31\_} Lemma 31.1). Refereed.}

\begin{proof}
Put $s_j=2\pi ij/\log a$ for $j\ge1$. For $n=a^k$, $n^{-s_j}=e^{-2\pi ijk}=1$, so every generator with $n$ a power of $a$ vanishes at every $s_j$.
Distinct point evaluations are linearly independent on $\B$ (the functions $e^{s^2}\prod_{i\ne j}(s-s_i)\in\B$ separate finitely many points), so among the evaluations at $s_1,\dots,s_{|F|+1}$ some nonzero combination $\Lambda=\sum_jc_j\operatorname{ev}_{s_j}$ satisfies the $|F|$ linear conditions $\Lambda(g_t(n-n^{1-s})/s)=0$, $n\in F$; it annihilates $\mathcal{N}_S$.
For the ideal, choose in the same way a nonzero combination $\Lambda'=\sum_jc'_j\operatorname{ev}_{s_j}$ annihilating $\xi g_t(n-n^{1-s})/s$ for $n\in F$; it annihilates $\xi\mathcal{N}_S$. It is nonzero on $I_\zeta$: $I_\zeta$ contains $g_t\xi P$ for every polynomial $P$, and $\Lambda'(g_t\xi P)=\sum_jc'_jg_t(s_j)\xi(s_j)P(s_j)\ne0$ for a suitable $P$, because $g_t(s_j)\ne0$ and $\xi(s_j)=\xi(1-s_j)\ne0$, $\zeta$ having no zeros on $\re s=1$.
\end{proof}

The programme's corrected tests behave like $\mathcal{N}_S$: if $\Phi$ is entire, $\Phi(0)=1$ and $g_t(1-\Phi)/s\in\B$, then $\{g_t(n^{1-s}-n\Phi)/s:n\ge1\}$ spans a dense subspace of $\B$ (register S59, \note{30\_} Corollary 30.2), and for the programme's correction $C=8F_*$ every tail $\{g_t(\varphi_m+C/s):m\ge M\}$ with Noor's $\varphi_m$ is already total in $\B$; for a general entire $C$ with $C(0)=1$ and $g_tC\in\B$, and $M\ge1$, the annihilator in $\B'$ of the tail is $0$ if $C(1)\ne0$ and $\C\operatorname{ev}_1$ if $C(1)=0$ (register S62, \note{32\_} Lemma 32.1).

\subsection{Positive forms and receivers}\label{ssec:receivers}

The programme examines many ways of attaching a positive form, a Hilbert space or a pairing to the quotient. The exact results are these.

\begin{lemma}[positive transfer-adjoint forms]\label{lem:PTQ}
On $H=\ell^2(\Zs,m)$ let $(T_nx)_\rho=n^\rho x_\rho$ and $U_n=nT_{1/n}$. A bounded positive semidefinite Hermitian form $B$ on $H$ satisfies $B(T_nx,y)=B(x,U_ny)$ for all integers $n\ge1$ if and only if $B(x,y)=\sum_{\re\rho=1/2}m_\rho c_\rho x_\rho\cl{y_\rho}$ with $0\le c_\rho\le C$.
\end{lemma}
\status{Proved in the programme (\lbl{PTQ4}); the off-diagonal step shortened here (claude-ab, \note{25\_} Lemma 25.2). Refereed (eighth pass).}

\begin{proof}
Since $U_nT_n=n$, the condition is $B(T_nx,T_ny)=nB(x,y)$. At $x=y=\varepsilon_\rho$ (the coordinate vector) it gives $(n^{2\re\rho}-n)B(\varepsilon_\rho,\varepsilon_\rho)=0$, so $B(\varepsilon_\rho,\varepsilon_\rho)=0$ off the line, and by Cauchy--Schwarz for $B$ such $\varepsilon_\rho$ pair to zero with everything.
For line zeros $\rho=\tfrac12+i\gamma$, $\eta=\tfrac12+i\gamma'$ it gives $(e^{i(\gamma-\gamma')\log n}-1)B(\varepsilon_\rho,\varepsilon_\eta)=0$. If $B(\varepsilon_\rho,\varepsilon_\eta)\ne0$ then $(\gamma-\gamma')\log2=2\pi k$ and $(\gamma-\gamma')\log3=2\pi l$; if $\gamma\ne\gamma'$ then $k\ne0$ and $3^k=2^l$, impossible. The converse is direct, since $|n^\rho|^2=n$ on the line.
\end{proof}

The following are proved in the programme and were audited; each is stated with its scope.
\begin{itemize}
\item \emph{Positive forms compatible with the scaling transfer} on $\B/I_\zeta$ are exactly $\sum_{\re\rho=1/2}m_\rho c_\rho F(\rho)\cl{G(\rho)}$ with polynomially bounded $c_\rho\ge0$ (\lbl{CFP}; \note{04\_}; register S12). Before the quotient, the positive transfer forms are exactly $\int F(\tfrac12+i\lambda)\cl{G(\tfrac12+i\lambda)}\,d\mu(\lambda)$ with $\mu$ positive and tempered, and they descend to the quotient if and only if $\operatorname{supp}\mu$ lies in the set of ordinates of critical zeros (\lbl{SPF0--SPF10}; \note{09\_}, \note{12\_}; S22; the Bochner--Schwartz theorem in Mellin coordinates).
\item \emph{All positive measure receivers at once.} For every positive measure $\mu$ with polynomially bounded mass on unit intervals, the closure of the image of $I_\zeta$ in $L^2(\mu)$ under $F\mapsto F(\tfrac12+i\lambda)$ is $L^2(\R\setminus\Gamma,\mu)$, $\Gamma$ the set of ordinates of critical zeros; so a positive receiver sees only line values at zeros carrying mass (\lbl{SMC1--SMC4}, \lbl{SMC9}; \note{18\_}; S35).
\item \emph{Connes against Meyer.} The natural map from $Q$ to the inverse limit of Connes' weighted Hilbert quotients \cite{Connes} has kernel exactly the classes vanishing on all critical-line jets, so it is injective if and only if RH holds; the algebraic map to each $H_\delta/J$ is injective, and the loss happens in the Hausdorff quotient (\lbl{WHR5--WHR10}; \note{12\_}; S28). The prime dilations on these Hilbert quotients have spectrum on $|z|=\sqrt p$ (\lbl{WHR8}; S27), and on the analogous quotients of $L^2(u^{2\sigma-1}(1+\log^2u)^\delta du)$ for $0<\sigma<1$ exactly the zeros on $\re s=\sigma$ are realised, with jet orders $j<\min(m_\rho,\delta-\tfrac12)$ and spectrum on $|z|=p^\sigma$; the family of these quotients over all $0<\sigma<1$ and all integers $\delta=N\ge1$ is jointly faithful on $Q$ (\lbl{VWR2--VWR7}; \note{14\_}; S31).
\item \emph{Spectral synthesis.} Explicit finite-rank Riesz--Gaussian operators converge to the identity on $Q$; closed multiplier submodules are exactly the jet-order submodules; the line ideal and the off-line ideal have dense sum (\lbl{GSP2--GSP6}, \lbl{CTS1--CTS4}; \note{12\_}; S29; read for correctness, not re-derived).
\item \emph{Global residue duality.} the pairing
\[
B_\zeta(F,G)=\frac1{2\pi i}\Big(\int_{\re s=2}-\int_{\re s=-1}\Big)\frac{F(s)G(1-s)}{\zeta(s)}\,ds
\]
satisfies $|B_\zeta(F,G)|\le\zeta(2)(1+4\pi^2)\pi^{-1}\,b_{2,1}(F)\,b_{2,1}(G)$, descends to $Q$, is nondegenerate, has local determinant $(m!/\zeta^{(m)}(\rho))^m$, and satisfies $B(T_ax,T_ay)=aB(x,y)$ and $B(Lx,y)=B(x,(1-L)y)$; it is not positive (\lbl{GZR}; \note{26\_} Lemma 26.1; S49).
\item \emph{The source-pairing identity.} If $b_0\in\mathcal{A}$ has Mellin transform vanishing at $\rho$ and $(L-\rho)b_\rho=b_0$ with $L=-u\partial_u$, then $(\re\rho-\tfrac12)\|b_\rho\|^2_{L^2(du)}=-\re\langle b_\rho,b_0\rangle$. (Proof: $\re\langle Lb,b\rangle=-\tfrac12\int u(|b|^2)'du=\tfrac12\|b\|^2$ by one integration by parts, so $\re\langle(L-\rho)b_\rho,b_\rho\rangle=(\tfrac12-\re\rho)\|b_\rho\|^2$.) It restates RH at $\rho$ as $\re\langle b_\rho,b_0\rangle=0$; reading it as a Hilbert--Pólya mechanism is an interpretation (\lbl{OPD4.2}; \note{09\_}; S20).
\item \emph{The harmonic defect and the divisor potential.} Poisson-sweeping the zeros onto the line gives $H(b_\dagger)=2\sum_\rho m_\rho|\re\rho-\tfrac12|\,\|b_\rho\|^2$ for an explicit test, and the swept distribution descends to $Q$ if and only if there are no off-line zeros (\lbl{HSW5--HSW6B}, \lbl{OPD5}; S21); for finitely many exponents with multiplicities, the translation trace compressed to the jet-constraint span on $[-T,T]$ tends to $\sum m\,e^{-|\beta-1/2||t|}e^{i\gamma t}$ (\lbl{FGR0--FGR6}; S26); and for $r>1$, $t>0$ the nonnegative quantity $\mathcal{E}_r(t)=\sum_\rho m_\rho(r^\sigma-r^{1-\sigma})^2e^{2t(\sigma^2-\gamma^2)}$, which vanishes if and only if RH holds, equals an area integral of the even part $\tfrac12\log|\zeta(s)\zeta(1-\bar s)|$ against the Laplacian of an explicit test function, plus $(r-1)^2(e^{2t}+1)/2$, when the test function is built from a symmetric cutoff equal to $1$ near $[0,1]$ and supported in $(-\tfrac12,\tfrac32)$; for a wider cutoff further terms from the trivial zeros appear (\lbl{OZD0--OZD7}; \note{12\_} §1.1; S25; a member of the Littlewood-lemma family \cite{BalazardSaiasYor,SekatskiiBeltraminelliMerlini}).
\end{itemize}
Each of these objects is exact, and none of them by itself decides where the zeros are: the positive ones see only critical zeros (a theorem, by the classifications above), the ones listed that see off-line zeros vanish, or are injective, exactly under RH, and the residue duality $B_\zeta$ is nondegenerate unconditionally and not positive. Section~\ref{sec:negative} records this as negative results with their scope.
